 % !TeX root = book.tex
\chapter{Implementing Web Prolog generics}\label{sec:implementing-prolog-generics}

Generic constructs provide reusable actor patterns that can be specialized,  configured or extended post creation-time, and that encapsulate common concerns such as message protocols, supervision and state management. %Conceptually, a generic behavior remains an actor in the usual sense; what distinguishes it is that the actor is wrapped in a deliberately engineered behavioural contract and runtime support system that enforces uniform structure and predictable semantics. An actor spawned directly with \texttt{spawn/1-3} is therefore a raw actor: an unconstrained thread of execution that must implement everything itself. A behavior, by contrast, structures and constrains the actor according to a predefined model of interaction.

It is useful to distinguish two broad categories of generics supported by Web Prolog. The first category comprises \emph{Erlang-style behaviors}, which are message-driven components exposing a stable protocol, typically realised as a fixed vocabulary of messages and replies. They do not participate directly in the caller's proof search or bind the caller's logic variables. Their purpose is to provide robust, reusable services with well-defined communication patterns. The  \texttt{server\_*} predicates implement a message-driven request-response service analogous to Erlang's \texttt{gen\_server} behavior, the \texttt{supervisor\_*} predicates provide a behavior similar to Erlang's supervisor behavior, and statechart actors implement an actor-based state machine model in the spirit of Erlang's \texttt{gen\_statem}. The \texttt{toplevel\_*} predicates also fall into this category, even though the abstraction they support -- a Prolog toplevel actor -- would not ordinarily be found in Erlang.

The second category comprises what may be termed \emph{Prolog-style generics}. Some of them are meta-predicates that internally orchestrate actors yet present a conventional predicate-level interface. From the programmer's perspective, they are ordinary predicates whose execution is implemented via actor orchestration behind the scenes. Examples include \texttt{parallel/1} and \texttt{first\_solution/2-3}, which spawn worker actors, coordinate them with selective receive, and enforce particular success, failure, and error propagation semantics. They are semi-deterministic and do not themselves enumerate solutions upon backtracking. By contrast, \texttt{rpc/2-3}, introduced in Chapter~4, is a Prolog-style behavior that can succeed nondeterministically with multiple solutions, since it acts as a distributed analogue of a predicate call and therefore participates directly in proof search.

This separation between Erlang-style behaviors and Prolog-style generics reflects a dual identity of the language. On one hand, Web Prolog inherits the actor model, including explicit message passing and robust supervision. On the other hand, it maintains Prolog's commitment to declarative interfaces and transparent variable binding. Behaviors provide a structured means of reconciling these two strands: they offer reusable actor frameworks where the machinery of communication, supervision, and concurrency is factored out from the application logic, and where programmers can work either at the level of actor protocols or at the level of ordinary predicates, depending on the needs of their application.

\section{The server behavior: a preliminary sketch}\label{sec:a-universal-stateful-server-with-hot-code-swapping}

The \emph{server behavior} implements a classic client-server pattern in Web Prolog: a central server process maintains a private state and serves an arbitrary number of clients that communicate through message passing. The goal is to obtain a generic, reusable, and stateful server that can be specialised by providing an application-specific callback predicate, while also supporting hot code swapping of that callback without interrupting service or losing state. This is a prototypical Erlang-style behavior: it exposes a stable message protocol without binding the caller's variables.

Here is a small convenience wrapper for spawning a server:

\begin{verbatim}
server_spawn(Pred, State, Pid) :-
    server_spawn(Pred, State, Pid, []).

server_spawn(Pred, State, Pid, Options) :-
    spawn(server_loop(Pred, State), Pid, Options).
\end{verbatim}

\noindent The predicate \texttt{server\_spawn/3-4} takes the name of a callback predicate \texttt{Pred/4}, an initial state \texttt{State}, and spawns a new server process running \texttt{server\_loop/2}. In this design, the state is always owned by the server process and never shared directly with clients. The behaviour of the server itself is captured by the following loop:

\begin{verbatim}
server_loop(Pred, State0) :-
    receive({
        '$call'(From, Ref, Request) ->
            call(Pred, Request, State0, Response, State),
            From ! Ref-Response,
            server_loop(Pred, State) ;
        '$upgrade'(Pred1) ->
            server_loop(Pred1, State0) ;
        '$stop'(From) ->
            From ! reply(true)
    }).
\end{verbatim}

\noindent When the pattern \texttt{'\$call'(From,\,Ref,\,Request)} matches an incoming message, \texttt{call(Pred,\,Request,\,State0,\,Response,\,State)} invokes the callback predicate \texttt{Pred/4}, where \texttt{State0} is the current server state and \texttt{State} is the updated state after processing the request. The server then replies to the client \texttt{From} with a message \texttt{Ref-Response} and recurs with the new state. All application logic lives in \texttt{Pred/4}; the generic server merely routes requests, manages the state parameter, and returns responses.

The second message form \texttt{'\$upgrade'(Pred1)} implements a very simple kind of hot code swapping. Upon receiving such a message, the server discards the old callback \texttt{Pred} and continues the loop with the new callback \texttt{Pred1}, while keeping the current state \texttt{State0} unchanged. In this way we can change the server's behaviour at runtime without stopping the process or reinitialising its state.

Here is the implementation of a suitable client API:

\begin{verbatim}
server_request(To, Request, Response) :-
    server_request(To, Request, Response, []).

server_request(To, Request, Response, Options) :-
    server_promise(To, Request, Ref),
    server_yield(Ref, Response, Options).
\end{verbatim}

\noindent The predicate \texttt{server\_request/3-4} behaves as a synchronous remote procedure call: it sends a request to the server and blocks until the matching reply is received or a timeout is triggered (depending on the \texttt{Options} passed to \texttt{receive/2} via \texttt{server\_yield/3}). Under the surface, it is implemented in two stages, separating the sending of a request from the collection of the corresponding response.

First, \texttt{server\_promise/3} sends the request and returns a fresh reference that will identify the response:

\begin{verbatim}
server_promise(To, Request, Ref) :-
    self(Self),
    make_ref(Ref),
    To ! '$call'(Self, Ref, Request).
\end{verbatim}

\noindent Here \texttt{make\_ref/1} creates a unique reference value, and the client process identifier \texttt{Self} is bundled with the request. Since the reference is sufficiently unique, multiple outstanding requests from the same client can safely be in flight at the same time; replies can arrive in any order and will still be matched correctly.

Second, \texttt{server\_yield/2-3} waits for the reply tagged with the given reference:

\begin{verbatim}
server_yield(Ref, Response) :-
    server_yield(Ref, Response, []).

server_yield(Ref, Response, Options) :-
    receive({
        Ref-Response -> true
    }, Options).
\end{verbatim}

\noindent The receive pattern \texttt{Ref-Response} ensures that only the response corresponding to the given reference is accepted; any other messages remain in the mailbox. In effect, this gives us a synchronous call interface on top of the underlying asynchronous message passing, with explicit correlation of calls and replies.

To specialise the generic server into a concrete service, we define an appropriate \texttt{Pred/4}. For example, in the refrigerator simulation, the callback \texttt{fridge/4} implements a tiny key-value store where the keys are food items and the value is a multiset of stored items:

\begin{verbatim}
fridge(store(Food), FoodList, ok, [Food|FoodList]).
fridge(take(Food), FoodList, ok(Food), Rest) :-
    select(Food, FoodList, Rest), !.
fridge(take(_Food), FoodList, not_found, FoodList).
\end{verbatim}

\noindent Using this definition we can start a server, store some food, and later retrieve it using the synchronous \texttt{server\_request/3}:

\begin{verbatim}
?- server_spawn(fridge, [], Pid).
Pid = 99491660.
?- server_request($Pid, store(milk), Response).
Response = ok.
?- server_request($Pid, take(milk), Response).
Response = ok(milk).
?-
\end{verbatim}

\noindent For applications where asynchronous interaction is preferred, here is how to use \texttt{server\_promise/3} and \texttt{server\_yield/2}:

\begin{verbatim}
?- server_promise($Pid, store(meat), Ref).
Ref = 11025019.
% ... perhaps do something else here ...
?- server_yield($Ref, Response).
Response = ok.
?-
\end{verbatim}

\noindent Upgrading a running server is just another message:

\begin{verbatim}
server_upgrade(To, Pred) :-
    To ! '$upgrade'(Pred).
\end{verbatim}

\noindent A client that knows the server's \texttt{Pid} (or registered name) can call \texttt{server\_upgrade/2} to replace the callback predicate. This operation is instantaneous from the client's perspective and does not wipe the server's state, which is exactly what we want from hot code swapping in a stateful service.

Now suppose we have developed a more efficient implementation \texttt{fridge2/4}, perhaps using a better data structure representing the state. As long as \texttt{fridge2/4} obeys the same interface \texttt{Request\,+\,OldState\,->\,Response\,+\,NewState}, we can upgrade the server in place:

\begin{verbatim}
?- server_upgrade(Pid, fridge2).
true.
?-
\end{verbatim}

\noindent The server process continues to run, all existing state (the current contents of the fridge) is preserved, and subsequent requests are handled by the new callback. This small example demonstrates the essence of the server behavior: a generic, reusable scaffold that provides state management, synchronous request-response communication, and lightweight hot code swapping, while leaving the application-specific logic to the callback predicate.

We terminate the server by calling \texttt{server\_stop/2}, which sends it a \texttt{'\$stop'} message:

\begin{verbatim}
server_stop(To, Reply) :-
    self(Self),
    To ! '$stop'(Self),
    receive({
        reply(Reply) -> true
    }).
\end{verbatim}


\subsubsection{Monitoring and fail-fast server calls}

The basic server protocol presented above supports synchronous request--response interaction by tagging each call with a fresh reference and waiting for a reply of the form \texttt{Ref-Response}. However, as written, a client may block indefinitely if the server terminates between receiving the request and sending the reply.

This is a natural place to use the \texttt{monitor/2} predicate introduced in Chapter~\ref{sec:web-prolog}. Before sending a synchronous request, the client installs a monitor on the server. The call then waits for \emph{either} the expected reply \texttt{Ref-Response} \emph{or} a \texttt{down} message indicating that the server has terminated. This yields a \emph{fail-fast} call abstraction: if the server dies, the waiting client is immediately notified and can raise an exception or otherwise propagate failure.

Unlike the \texttt{monitor(true)} option, also introduced in Chapter~\ref{sec:web-prolog}, \texttt{monitor/2} is not atomic with process creation: it is a separate operation performed \emph{after} the target process already exists. This means that a server could, in principle, terminate in the small gap between obtaining its pid and installing the monitor. Nevertheless, \texttt{monitor/2} remains essential in practice, precisely because it applies to actors that already exist -- for example servers discovered through registration, servers started long before the current client, or services whose pid is obtained from another component. In such cases, spawn-time monitoring is simply not available, and \texttt{monitor/2} is the appropriate tool.

We implement fail-fast calls by extending \texttt{server\_promise/3} to also return a monitor reference, and by extending \texttt{server\_yield/2-3} to accept either a reply or the corresponding \texttt{down} message:

\begin{verbatim}
server_promise(To, Request, Ref, MonRef) :-
    self(Self),
    monitor(To, MonRef),
    make_ref(Ref),
    To ! '$call'(Self, Ref, Request).

server_yield(Ref, MonRef, Response, Options) :-
    receive({
        Ref-Response0 ->
            demonitor(MonRef),
            Response = Response0 ;
        down(MonRef, _Pid, Reason) ->
            throw(server_down(Reason))
    }, Options).
\end{verbatim}

\noindent Consequently, \texttt{server\_request/4} must then be implemented like this:

\begin{verbatim}
server_request(To, Request, Response, Options) :-
    server_promise(To, Request, Ref, MonRef),
    server_yield(Ref, MonRef, Response, Options).
\end{verbatim}

\noindent The monitor reference \texttt{MonRef} is essential: it identifies which monitor instance produced the \texttt{down} message, and it allows the client to cancel the monitor once the reply has been received.

Here is what happens when something goes wrong:

\begin{verbatim}
?- server_spawn(fridge, [], Pid, [monitor(true)]).
Pid = 41168156.
?- server_request($Pid, sore(milk), Response).
ERROR: Unhandled exception: Unknown message: server_down(false)
?- flush.
Shell got down(41168156,41168156,false)
true.
?-
\end{verbatim}

\noindent The server crashed because we accidentally requested \texttt{sore(milk)} rather than \texttt{store(milk)}. The current version of \texttt{fridge/4} could not handle this.

\medskip

\noindent The server behaviour illustrates a principle that runs through the rest of
this chapter and, more broadly, through the design of the Prolog Trinity as
a whole. What we have done is cleanly separate the fixed scaffolding of a
stateful concurrent service -- its message protocol, state threading,
monitoring, and support for code upgrade -- from the application-specific
logic, which is captured entirely in the callback predicate. The developer
writes only the part that expresses the domain semantics; the generic server
handles all the operational machinery.

It is worth pausing to compare this pattern to the simple refrigerator
simulation in Chapter~\ref{sec:web-prolog}. There, we modelled the behaviour of a fridge
directly as a Prolog relation over states. Here, we take exactly the same
relational specification and plug it into a reusable server framework that
gives it a process identity, encapsulated state, synchronous or
asynchronous interaction, and hot code swapping. Nothing about the original
Prolog style is lost. Developers can still write ordinary pure predicates to
describe how a state evolves, but they now have the option of turning such
predicates into fully fledged concurrent services with minimal additional
code. This is the distinctive promise of the approach: it elevates familiar
Prolog programming into a concurrent and distributed setting without
requiring a change in declarative style.

%\noindent So what was the win here? We have cleanly separated the fixed scaffolding of a stateful server -- its message protocol, state threading, and support for code upgrade -- from the application-specific logic, which is captured entirely in the callback predicate. In other words, the developer writes \emph{only} the part that expresses the domain semantics, while the generic server handles all the concurrent and operational machinery. It is worth comparing this pattern to the simple refrigerator simulation in Chapter~2. There, we modelled the behaviour of a fridge directly as a Prolog relation over states. Here, we take exactly the same relational specification and ``plug it into'' a reusable server framework that gives it a process identity, encapsulated state, synchronous or asynchronous interaction, and hot code swapping. Nothing about the original Prolog style is lost: developers can still write ordinary pure predicates to describe how a state evolves, but they now have the option of turning such predicates into fully fledged concurrent services with minimal additional code. This is the main advantage of the server behaviour: it elevates familiar Prolog programming into a concurrent and distributed setting without requiring a change in declarative style.

Of course, this preliminary sketch omits some of the practical details that a production-quality server behavior would need to address. Nevertheless, it already captures the core pattern that might be refined and extended later.

\section{The supervisor behavior}

Fault tolerance is essential in software applications because failures are
inevitable in real systems, and a fault-tolerant design ensures that
individual errors do not cascade into system-wide outages but are
contained, recovered from, and rendered harmless to users and neighbouring
components.

In Erlang, fault tolerance is achieved through supervision trees~--
hierarchies of actors that contain and isolate crashes. At the core of this
model are \emph{supervisors}, actors whose sole purpose is to start child
actors, monitor them, and decide how to recover from crashes. Supervisors
themselves can be supervised, yielding supervision trees of, in principle,
arbitrary depth.

Underpinning this approach is the so-called \emph{let it crash} principle:
worker actors are not burdened with complex error handling but are allowed
to crash when something goes wrong. The supervisor detects the crash
through a \texttt{down} message containing a pid, identifies the crashed
actor by this pid, and applies a predefined restart strategy. Erlang's
restart policies are uniform and predictable. Depending on the severity and
scope of the crash, a supervisor may do nothing, restart the crashed actor
process (and perhaps its siblings), or terminate itself so that its own
supervisor can take corrective action.

The implementation that follows is deliberately minimal: a one-for-one
restart policy with a single global restart budget, enough to show how the
actor primitives compose into a supervision structure but far short of a
production-quality supervisor. The richer mechanisms that a full
implementation would require~-- alternative restart strategies, per-child
restart intensities, dynamic child management~-- are discussed at the end
of this section. Our purpose here is not completeness but demonstration: to
confirm that the same \texttt{spawn}, \texttt{monitor}, \texttt{exit}, and
selective \texttt{receive} primitives that built the server behaviour in
the previous section are sufficient to support fault-tolerant process
management as well.

A supervisor consists of a \emph{specific part}, which defines which child
processes to manage and under what policies they are started and restarted,
and a \emph{generic part}~-- responsible for starting, linking, monitoring,
and restarting children. In our implementation, the specific part is given
by a list of child specifications and the generic part is captured by the
supervisor loop. Each child specification carries an identifier, a start
goal, and a restart policy: \emph{permanent} children are always restarted,
\emph{transient} children are restarted only if they terminate abnormally,
and \emph{temporary} children are never restarted.

The architectural heart of the supervisor is its main loop:
%
\begin{verbatim}
supervisor_loop(ChildList, Count0) :-
    receive({
        down(_, Pid, true) ->
            remove_child(Pid, ChildList, ChildList1),
            supervisor_loop(ChildList1, Count0) ;
        down(_, Pid, Why) ->
            (   Count0 > 0
            ->  Count is Count0 - 1,
                restart_child(Pid, Why, ChildList, ChildList1),
                supervisor_loop(ChildList1, Count)
            ;   exit(gave_up)
            ) ;
        '$stop'(From) ->
            terminate(ChildList),
            From ! reply(true)
    }).
\end{verbatim}
%
The loop waits in a \texttt{receive} call for \texttt{down} messages from
monitored children. If the exit reason is \texttt{true} (normal
termination), the child is simply removed from the child list. If the
reason indicates a crash, and the restart budget has not been exhausted, the
crashed child is restarted and its entry in the child list replaced with
the new pid. If the budget is exhausted, the supervisor terminates
itself~-- which, if it is itself supervised, allows a higher-level
supervisor to take corrective action. Upon receiving a stop message, the
supervisor shuts down all children in an orderly fashion.

Here is a brief example. Given three workers~-- one that blocks forever,
one that terminates normally after three seconds, and one that crashes
after five~-- we start a supervisor with a budget of three restarts:\footnote{The option and child-spec forms used in this sketch (notably \texttt{budget/1} and \texttt{child/3}) are provisional. For the normative, implemented supervisor API, see Appendix~\ref{sec:manual}.}
%
\begin{verbatim}
?- supervisor_spawn([
       child(steady,  (receive({})),   permanent),
       child(runner,  sleep(3),        temporary),
       child(crasher, (sleep(5), exit(an_error)), permanent)
   ], Pid, [budget(3), name(my_supervisor), monitor(true)]).
Pid = 10039133.   
?-
\end{verbatim}
%
The crasher terminates abnormally and is restarted. After three restart
attempts the budget is exhausted and the supervisor gives up, terminating
itself. The remaining child (the steady worker) is killed by the link
cascade when the supervisor dies~-- no orphaned processes are left behind.

This sketch omits several mechanisms found in a production-quality
supervisor. Erlang's OTP, for example, supports configurable restart
strategies~-- \texttt{one\_for\_all} and \texttt{rest\_for\_one}, which
dictate whether sibling processes must also terminate when a crash
occurs~-- as well as per-child restart intensities with time windows (e.g.\
three crashes in ten seconds), and a robust API for dynamically adding or
removing child specifications at runtime. None of these extensions would
require new primitives. Each is implementable using the same
\texttt{spawn}, \texttt{monitor}, \texttt{exit}, and selective
\texttt{receive} that the sketch already relies on, and in each case the
callback-based structure would remain: the supervisor's generic loop
handles lifecycle management while the child specifications capture the
application-specific topology. A complete implementation covering all three
restart strategies, dynamic child management, restart intensity tracking,
nested supervision, and cascade cleanup~-- together with a test suite of
sixteen passing tests~-- is available on the book's companion
page.\footnote{\url{https://...}} % TODO: insert URL
The point is not that this sketch is complete, but that the actor substrate
introduced in Chapter~2 is sufficient to support supervision patterns of
arbitrary sophistication. For a detailed account of OTP's full supervision
capabilities, see the Erlang
documentation.\footnote{\url{https://www.erlang.org/doc/system/sup\_princ.html}}


%\section{The supervisor behavior: a preliminary sketch}\label{sec:supervisor-example}
%
%\textit{Fault tolerance} is essential in software applications because failures are inevitable in real systems, and a fault-tolerant design ensures that individual errors do not cascade into system-wide outages but are contained, recovered from, and rendered harmless to users and neighbouring components.
%
%In Erlang, fault tolerance is achieved through \emph{supervision trees} -- hierarchies of actors that contain and isolate crashes. At the core of this model are \emph{supervisors}, actors whose sole purpose is to start child actors, monitor them, and decide how to recover from crashes. Supervisors themselves can be supervised, yielding supervision trees of, in principle, arbitrary depth.
%
%Underpinning this approach is the so-called \emph{let it crash} principle: worker actors are not burdened with complex error handling but are allowed to crash when something goes wrong. The supervisor detects the crash through a message containing a pid, identifies the crashed actor by this pid, and applies a predefined restart strategy.
%
%Erlang's restart policies are uniform and predictable. Depending on the severity and scope of the crash, a supervisor may do nothing, restart the crashed actor process (and perhaps its siblings), or terminate itself so that its own supervisor (if any) can take corrective action. This uniformity makes the system's behavior reliable and comprehensible, even in complex crash scenarios.
%
%The implementation that follows is deliberately minimal: a one-for-one
%restart policy with a single global restart budget, enough to show how the
%actor primitives compose into a supervision structure but far short of a
%production-quality supervisor. The richer mechanisms that a full
%implementation would require -- alternative restart strategies, per-child
%restart intensities, dynamic child management -- are enumerated at the end
%of this section. Our purpose here is not completeness but demonstration: to
%confirm that the same \texttt{spawn}, \texttt{monitor}, \texttt{exit}, and
%selective \texttt{receive} primitives that built the server behaviour in the
%previous section are sufficient to support fault-tolerant process management
%as well.
%
%A supervisor consists of a \emph{specific} part, which defines which child processes to manage and under what policies they are started and restarted, and a \emph{generic} part -- responsible for starting, linking, monitoring, and restarting children. In the terminology of this section, a supervisor is an Erlang-style behavior specialised by a set of child specifications.
%
%We proceed to sketch a simple supervisor in Web Prolog. In our implementation, the specific part is given by a list of \emph{child specifications}. The only restart policy supported by the generic part is straightforward: a process that terminates normally is allowed to come to rest, whereas a process that crashes will be restarted -- but, as a safeguard against endless cycles of failure, only up to a limited number of times.
%
%Our supervisor is started by calling \texttt{supervisor\_spawn/2-3} with a list of child specifications. Two options are supported: \texttt{name} is for giving the supervisor a registered name, and the option \texttt{budget} takes an integer that determines how many times a restart is performed before the supervisor gives up:
%
%\begin{verbatim}
%supervisor_spawn(ChildSpecList, Pid) :-
%    supervisor_spawn(ChildSpecList, Pid, []).
%
%supervisor_spawn(ChildSpecList, Pid, Options) :-
%    option(budget(Count), Options, 5),
%    spawn(initialize(ChildSpecList, Count), Pid, Options),
%    (   option(name(Name), Options)
%    ->  register(Name, Pid)
%    ;   true
%    ).
%\end{verbatim}
%
%\noindent The predicate \texttt{initialize/2} spawns all of the child processes that will be supervised and creates a list of pairs of the form \texttt{Pid-ChildSpec} that is then passed to the \texttt{supervisor\_loop/2} predicate:
%
%\begin{verbatim}
%initialize(ChildSpecList, Count) :-
%    maplist(start_child, ChildSpecList, ChildList),
%    supervisor_loop(ChildList, Count).
%
%start_child(child(ID, Spec, Restart), Child) :-
%    do_start(Spec, ID, Pid),
%    Child = Pid-child(ID, Spec, Restart),
%    format("Child ~w started with pid ~w~n", [ID, Pid]).
%\end{verbatim}
%
%\noindent Here is how to start different kinds of children:
%
%\begin{verbatim} 
%do_start(server(Pred, ServerOpts), Name, Pid) :- !,
%    append(ServerOpts, [name(Name), monitor(true)], Opts1),
%    server_spawn(Pred, Pid, Opts1).
%do_start(Goal, _Name, Pid) :-
%    spawn(Goal, Pid, [monitor(true)]).
%\end{verbatim}
%
%\noindent The supervisor loop waits in a receive call for \texttt{down} and \texttt{'\$stop'} messages. If a child terminates, the supervisor receives a message of the form \texttt{down(Ref,\,Pid,\,Reason)}. If \texttt{Reason} is bound to \texttt{true} the child is removed from the child list. If \texttt{Count0\,>\,0}, then the crashed child is restarted with a new pid and its entry in the child list replaced, or else the supervisor gives up:
%
%\begin{verbatim}
%supervisor_loop(ChildList, Count0) :-
%    receive({
%        down(_, Pid, true) ->
%            format("~w halted normally~n", [Pid]),
%            remove_child(Pid, ChildList, ChildList1),
%            supervisor_loop(ChildList1, Count0) ;
%        down(_, Pid, Why) ->
%            format("~w crashed because of ~w~n", [Pid, Why]),
%            (   Count0 > 0
%            ->  Count is Count0 - 1,
%                restart_child(Pid, Why, ChildList, ChildList1),
%                supervisor_loop(ChildList1, Count)
%            ;   format("Gave up on ~w, and halted~n", [Pid]),
%                exit(gave_up)
%            ) ;
%        '$stop'(From) ->
%            terminate(ChildList),
%            From ! reply(true)
%    }).
%
%remove_child(Pid, ChildList0, ChildList) :-
%    once(select(Pid-_, ChildList0, ChildList)).
%
%restart_child(Pid, Why, ChildList, NewChildList) :-
%    once(select(Pid-child(ID, Spec, Restart), ChildList, Rest)),
%    (   restart_allowed(Restart, Why)
%    ->  do_start(Spec, ID, NewPid),
%        format("~w restarted with pid ~w~n", [ID, NewPid]),
%        NewChildList = [NewPid-child(ID, Spec, Restart) | Rest]
%    ;   format("No restart: ~w (policy ~w)~n",[ID, Restart]),
%        NewChildList = Rest
%    ).
%\end{verbatim}
%
%\noindent The predicate \texttt{restart\_allowed/2} defines when a terminated child process must be restarted. With terminology borrowed from Erlang, a \texttt{permanent} child process is always restarted. A \texttt{temporary} child process is never restarted. A \texttt{transient} child process is restarted only if it terminates abnormally, that is, with another exit reason than \texttt{true}:
%
%\begin{verbatim}  
%restart_allowed(temporary, _Why) :- !, fail.
%restart_allowed(transient, true)    :- !, fail.
%restart_allowed(_Restart, _Why).            
%\end{verbatim}
%
%\noindent Upon receiving the \texttt{stop} message, the supervisor runs through the child list, terminating the children one by one:
%
%\begin{verbatim}     
%terminate([]).        
%terminate([Pid-_|ChildList]) :- 
%    exit(Pid, kill),
%    terminate(ChildList).
%\end{verbatim}
%
%\noindent We stop the supervisor by calling \texttt{supervisor\_stop/2}, which can be written in the same style as \texttt{server\_stop/2}, so no need to show its implementation here.
%
%As in Erlang, processes must be started in a particular manner to be eligible for supervision. Our supervisor only requires that child specifications are callable goals that can be spawned. In preparation for a trial run, the following trivial definitions specify three worker processes: the first runs forever, the second sleeps for three seconds and then terminates normally, the third process sleeps for five seconds and then crashes with an error:
%
%\begin{verbatim}
%actor_1 :- receive({}).    
%\end{verbatim}
%\vspace{-0.3cm}
%\begin{verbatim}
%actor_2 :- sleep(3).    
%\end{verbatim}
%\vspace{-0.3cm}
%\begin{verbatim}
%actor_3 :- sleep(5), exit(an_error).
%\end{verbatim}
%
%\noindent Here is an example system startup procedure:
%
%\begin{verbatim}
%demo_supervisor :-
%    ChildSpecs = [
%        child(fridge_server,
%              server(fridge, [initial_state([])]),
%              permanent),
%        child(worker1, actor_1, transient),
%        child(worker2, actor_2, temporary),
%        child(worker3, actor_3, permanent)  
%    ],
%    supervisor_spawn(ChildSpecs, Pid, [
%        budget(3), 
%        name(my_supervisor),
%        monitor(true)
%    ]),
%    format("Demo started with pid ~w~n", [Pid]),
%    receive({
%        down(_, Pid, Why) ->
%            format("Demo halted with reason ~w~n", [Why])
%    }).
%\end{verbatim}
%
%\noindent Since the code is instrumented with calls to \texttt{format/2}, it is easy to follow what it does. Here is a trial run that shows that the supervisor works as expected:
%
%\begin{verbatim}
%?- demo_supervisor.
%Demo started with pid 25907375
%fridge_server started with pid 36917416
%worker1 started with pid 21857410
%worker2 started with pid 44503190
%worker3 started with pid 98783705
%44503190 terminated normally
%98783705 crashed because of an_error
%worker3 restarted with pid 26352791
%26352791 crashed because of an_error
%worker3 restarted with pid 91138479
%91138479 crashed because of an_error
%worker3 restarted with pid 43153221
%43153221 crashed because of an_error
%Gave up on 43153221, and halted
%Demo halted with reason gave_up
%true.
%?-
%\end{verbatim}
%
%\noindent Recall that \texttt{spawn/2-3} links children to their parent process by default, so when the supervisor terminates abnormally, the  only remaining child actor (running \texttt{actor\_1/0}) was forced to terminate too.
%
%This sketch omits several mechanisms found in a production-quality
%supervisor. Erlang's OTP, for example, supports configurable restart
%strategies -- \texttt{one\_for\_all} and \texttt{rest\_for\_one}, which
%dictate whether sibling processes must also terminate when a crash occurs --
%as well as per-child restart intensities with time windows (e.g.\ three
%crashes in ten seconds), and a robust API for dynamically adding or removing
%child specifications at runtime. None of these extensions would require new
%primitives. Each is implementable using the same \texttt{spawn},
%\texttt{monitor}, \texttt{exit}, and selective \texttt{receive} that the
%sketch already relies on, and in each case the callback-based structure
%would remain: the supervisor's generic loop handles lifecycle management
%while the child specifications capture the application-specific topology.
%The point is not that this sketch is complete, but that the actor substrate
%introduced in Chapter~\ref{sec:web-prolog} is sufficient to support supervision patterns of
%arbitrary sophistication. For a detailed account of OTP's full supervision
%capabilities, see the Erlang
%documentation.\footnote{\url{https://www.erlang.org/doc/system/sup\_princ.html}}


%While functional, this lightweight implementation omits several sophisticated reliability mechanisms found in Erlang's Open Telecom Platform (OTP). Standard OTP supervisors support configurable restart strategies -- such as \texttt{one\_for\_all} or \texttt{rest\_for\_one} -- which dictate whether sibling processes must also terminate when a crash occurs. Furthermore, our example uses a single global restart counter for all children, instead of per-child restart intensities and time windows (e.g., 3 crashes in 10 seconds). It also lacks a robust API for dynamically adding or removing child specifications at runtime without stopping the supervision loop. However, the core logic presented here forms a solid foundation; by extending our example considerably to include these features, a sophisticated supervisor fully capable of managing complex process trees can be built. For a detailed account of these capabilities, see the Erlang documentation.\footnote{\url{https://www.erlang.org/doc/system/sup\_princ.html}}

\section{The toplevel behavior}\label{sec:toplevel-implementation}

The toplevel behavior, controlled by predicates such as \texttt{toplevel\_spawn/1-2} and friends, must also be implemented. In our experience, once we have a complete implementation of all the Erlang-style primitives for concurrency and distribution, the implementation of the toplevel behavior and the built-in predicates for controlling it is fairly straightforward.

We begin with an implementation of \texttt{toplevel\_spawn/1-2}:

\begin{verbatim}
toplevel_spawn(Pid) :-
    toplevel_spawn(Pid, []).

toplevel_spawn(Pid, Options) :-
    self(Self),
    option(target(Target), Options, Self),
    option(session(Continue), Options, false),
    spawn(session(Pid, Target, Continue), Pid, Options).
\end{verbatim}

\noindent Before we proceed to implement \texttt{session/3} we shall demonstrate a couple of predicates that would make the task much easier. SWI-Prolog offers a buil-in predicate \texttt{findnsols/4} which provides a useful foundation for our implementation. It is somewhat similar to the standard \texttt{findall/3}, but expects an integer \texttt{Limit} in its first argument and will generate at most that many solutions. It is also nondeterministic, so on backtracking it will do it again. We borrow an example of its use from the SWI-Prolog manual:\footnote{\url{https://www.swi-prolog.org/pldoc/doc\_for?object=findnsols/4}}

\begin{verbatim}
?- findnsols(5, I, between(1, 12, I), L).
L = [1, 2, 3, 4, 5] ;
L = [6, 7, 8, 9, 10] ;
L = [11, 12].
?-
\end{verbatim}

\noindent Another SWI-Prolog library predicate \texttt{offset/2} will also prove useful.\footnote{https://www.swi-prolog.org/pldoc/doc\_for?object=offset/2} Its purpose is to \textit{skip} the first \textit{n} solutions to a goal, i.e. the first \textit{n} solutions are computed, but not collected. Here is an example of its use:

\begin{verbatim}
?- offset(10, between(1, 12, I)).
I = 11 ;
I = 12.
?- 
\end{verbatim}

\noindent Combining \texttt{findnsols/4} with \texttt{offset/2} allows us to implement a predicate \texttt{slice/5} capable of computing a \textit{slice} of solutions to a goal:

\begin{verbatim}
slice(Goal, Template, Offset, Limit, Slice) :-
   findnsols(Limit, Template, offset(Offset, Goal), Slice).
\end{verbatim}

\noindent However, we are looking for \textit{answers}, rather than just slices of solutions. By wrapping a call to \texttt{slice/5} in a call to \texttt{call\_cleanup/2} wrapped by a call to \texttt{catch/3} we arrive at a predicate \texttt{answer/5} capable of producing the four different forms of answer terms that we need:

\begin{verbatim}
answer(Goal, Template, Offset, Limit, Answer) :-
   catch(
      call_cleanup(slice(Goal, Template, Offset, Limit, Slice),
                   Det = true),
         Error, true),
   (   Slice == []
   ->  Answer = failure
   ;   nonvar(Error)
   ->  Answer = error(Error)
   ;   var(Det)
   ->  Answer = success(Slice, true)
   ;   Det = true
   ->  Answer = success(Slice, false)
   ).
\end{verbatim}


\noindent Moving along to the implementation of \texttt{session/3} and the PTCP, consider again the statechart in Figure~\ref{fig:complate-states-generic}:

\begin{figure}[h]
    \centering
	\includegraphics[width=9cm]{PLAP-statechart-3-states.pdf}
    \caption{The complete Prolog Toplevel Communication Protocol.}
    \label{fig:complate-states-generic}
\end{figure}

\noindent The most important part of the implementation of the PTCP are the three inner states \textbf{s1}, \textbf{s2} and \textbf{s3}.
We shall ignore for a moment the outer state labelled PTCP and implement a session like so:

\begin{verbatim}
session(Pid, Target0, Continue) :-
    state_1(Pid, Target0, Continue).
\end{verbatim}

\noindent This takes us straight to \textbf{s1}, and here is what happens in the predicate \texttt{state\_1}, which serves as the implementation of this state:

\begin{verbatim}
state_1(Pid, Target0, Continue) :-
    receive({
        '$call'(Goal, Options) ->
            option(template(Template), Options, Goal),
            option(offset(Offset), Options, 0),
            option(limit(Limit0), Options, 10 000 000 000),
            option(target(Target1), Options, Target0),
            Limit = count(Limit0),
            state_2(Goal, Template, Offset, Limit, Pid, Answer),
            Target = target(Target1),
            arg(1, Target, Out),
            Out ! Answer,
            (   arg(3, Answer, true)
            ->  state_3(Limit, Target)  
            ;   true
            ) 
        }),
    (   Continue == false
    ->  true
    ;   state_1(Pid, Target0, Continue)
    ).
\end{verbatim}

\noindent In \texttt{state\_2} the real work is being done. This is where \texttt{answer/5}, defined earlier in this chapter, is used. However, answer terms must be extended with the pids of the actor processes that produced them. 

\begin{verbatim}
state_2(Goal, Template, Offset, Limit, Pid, Answer) :-
    answer(Goal, Template, Offset, Limit, Answer0),
    add_pid(Answer0, Pid, Answer).
\end{verbatim}

\noindent To handle this, \texttt{add\_pid/3} is defined like so:

\begin{verbatim}
add_pid(success(Slice, More), Pid, success(Pid, Slice, More)).
add_pid(failure, Pid, failure(Pid)).
add_pid(error(Term), Pid, error(Pid, Term)).
\end{verbatim}

\noindent One feature of \texttt{answer/5} that was not demonstrated before, is that the argument specifying the limit can be passed a term \texttt{count/1} with an integer in its argument. This works like a mutable local variable that can be assigned values using \texttt{nb\_setarg/3} and read by means of \texttt{arg/3}.

\begin{verbatim}
?- Limit = count(2),
   answer(between(1,12,I), I, 0, Limit, Answer),
   nb_setarg(1, Limit, 5).
Limit = count(5),
Answer = success([1, 2], true) ;
Limit = count(5),
Answer = success([3, 4, 5, 6, 7], true) ;
Limit = count(5),
Answer = success([8, 9, 10, 11, 12], false).
?-
\end{verbatim}

\noindent In the definition of the predicate \texttt{state\_1/3} we saw that if a \texttt{success} answer term indicates (with \texttt{true} in its third argument) that there may be more solutions to the current goal, we enter \texttt{state\_3}. For other answer terms a recursive call of \texttt{state\_1/3} is made.

\begin{verbatim}
state_3(Limit, Target) :-
    receive({
        '$next'(Options2) ->
            (   option(limit(NewLimit), Options2)
            ->  nb_setarg(1, Limit, NewLimit)
            ;   true
            ),
            (   option(target(NewTarget), Options2)
            ->  nb_setarg(1, Target, NewTarget)
            ;   true
            ),                                                         
            fail ;
        '$stop' -> true
    }).
\end{verbatim}

\noindent Here it is the reception of the \texttt{'\$next'} message and the subsequent call to \texttt{fail/0} that triggers the backtracking to \texttt{answer/5} in state s2. If the \texttt{'\$stop'} message is received instead, \texttt{state\_3/2} terminates, and then \texttt{state\_1/3} terminates too (unless the option \texttt{session(true)} was passed to \texttt{toplevel\_spawn/1-2}).

As can be seen in the diagram depicting the PTCP, we have so far only implemented the three inner states. 

We need to enable a client to abort the execution of a goal. We can do this by changing the implemention of \texttt{session/3} like so::

\begin{verbatim}    
session(Pid, Target, Continue) :-
    catch(state_1(Pid, Target, Continue), 
          '$abort_goal',
          session(Pid, Target, Continue)).
\end{verbatim}

\noindent Here is how to tell the toplevel actor to abort the execution of any goal that it currently runs:

\begin{verbatim}
toplevel_abort(Pid) :-
    catch(thread_signal(Pid, throw('$abort_goal')), 
          error(existence_error(_,_), _), 
          true).  
\end{verbatim}

\noindent The action of aborting a particular execution of a goal passed to \texttt{toplevel\_call/2-3} must not be confused with the action of exiting the toplevel process. The latter can be performed by calling \texttt{exit/1-2}.

Below, we implement \texttt{toplevel\_call/2-3} simply as

\begin{verbatim}
toplevel_call(Pid, Goal) :-
    toplevel_call(Pid, Goal, []).

toplevel_call(Pid, Goal, Options) :-
    Pid ! '$call'(Goal, Options).
\end{verbatim}

\noindent and \texttt{toplevel\_next/1-2} like so

\begin{verbatim}
toplevel_next(Pid) :-
    toplevel_next(Pid, []).

toplevel_next(Pid, Options) :-
    Pid ! '$next'(Options). 
\end{verbatim}    

\noindent and, finally, \texttt{toplevel\_stop/1} like so:

\begin{verbatim}
toplevel_stop(Pid) :-
    Pid ! '$stop'.
\end{verbatim}    

\section{The parallel/1 meta-predicate}

In this section we will implement \texttt{parallel/1}, a meta-predicate that tries to run a list of goals in parallel. A call to \texttt{parallel(Goals)} should

\begin{enumerate}
\item block until all work has been done, but no longer,
\item succeed, with variable bindings, if all goals succeed,
\item fail as quickly as possible if any goal fails, and 
\item rethrow any errors thrown by a goal, also as quickly as possible.
\end{enumerate}

\noindent In other words, running a list of goals in parallel should behave in the same way as when running them sequentially, but finish faster. For this to work properly, we must require something about the input too, namely that goals in the list are independent, that is, they must not communicate using shared variables or by any other means. From the caller's perspective, \texttt{parallel/1} is a Prolog-style behavior: a pure-looking meta-predicate whose implementation happens to orchestrate actors.

To facilitate the explanation of our approach to implementing \texttt{parallel/1}, we first define a predicate \texttt{parallel/2} that accepts exactly two goals and spawns two actor processes to solve them in parallel:

\begin{verbatim}
parallel(Goal1, Goal2) :-
    self(Self),
    spawn((call(Goal1), Self ! Pid1-Goal1), Pid1),
    spawn((call(Goal2), Self ! Pid2-Goal2), Pid2),        
    receive({Pid1-Goal1 -> true}),
    receive({Pid2-Goal2 -> true}).    
\end{verbatim}

\noindent The actor \texttt{Pid1} calls \texttt{Goal1} and, if it succeeds, sends the pair of \texttt{Pid1} and the (now instantiated) goal term as a message to the parent \texttt{Self}. The actor \texttt{Pid2} does the same thing for \texttt{Goal2}.

If the message sent by \texttt{Pid1} reaches the parent's mailbox first, then the first \texttt{receive} clause is triggered, and execution steps to the second \texttt{receive} statement and waits for the second actor's message to arrive. If the message sent by \texttt{Pid2} arrives first, then it is deferred. Once the message from \texttt{Pid1} comes along, the first \texttt{receive} statement is satisfied and the program steps to the second \texttt{receive} statement, which is triggered immediately by the deferred message. The result is that when \texttt{parallel/2} terminates, the messages will have been received irrespective of the order in which they were sent. The time spent waiting for the call to succeed is the longer of the response times from the two actors. As so often, asynchronous communication using send in combination with selective receive is key to the kind of programming going on here.

It is important to remember that the goal in the first argument of \texttt{spawn/1-3} is always copied before being called. This means that the instantiation of a goal passed to \texttt{parallel/2} will not happen until the corresponding call to \texttt{receive/1} has selected the message sent from that call.

This program satisfies our requirements 1) and 2), but not 3) and 4). The problem is that if one goal fails or throws an error, the corresponding \texttt{receive} will suspend, waiting in vain for a message of the form \texttt{Pid-Goal} to show up. We will explain how to deal with this, but first show how our approach can be generalised into taking a list of goals instead of just two. For this, \texttt{maplist/3} comes in handy: 

\begin{verbatim}
parallel(Goals) :-
    maplist(par_solve, Goals, Pids),
    maplist(par_yield, Pids, Goals).

par_solve(Goal, Pid) :- 
    self(Self),
    spawn((call(Goal), Self ! Pid-Goal), Pid).

par_yield(Pid, Goal) :-        
    receive({Pid-Goal -> true}). 
\end{verbatim}

\noindent This, by itself, does not help us satisfy the requirement 3) and 4), but here is a version of the program that does:

\begin{verbatim}    
parallel(Goals) :-
    maplist(par_solve, Goals, Pids),
    maplist(par_yield(Pids), Pids, Goals).

par_solve(Goal, Pid) :-
    self(Self),
    spawn((call(Goal), Self ! Pid-Goal), Pid, [
        monitor(true)
    ]).

par_yield(Pids, Pid, Goal) :-
    receive({
        Pid-Goal ->
            receive({
                down(_, Pid, true) ->
                    true
            }) ;
        down(_, _, false) ->
            tidy_up_all(Pids),
            !, fail ;
        down(_, _, exception(E)) ->
            tidy_up_all(Pids),
            throw(E)
    }).
\end{verbatim}

\noindent In this version, all three predicates from the previous program have been modified. In \texttt{par\_solve/2} each actor running a call is now being monitored, and \texttt{par\_yield/3} (which is \texttt{par\_yield/2} with the addition of a third argument, the purpose of which we shall explain in a bit) is set up to inspect the \texttt{down} messages of the form \texttt{down(Ref, Pid, Reason)} arriving from the worker actors and act appropriately. If \texttt{Reason} is \texttt{true} nothing needs to be done, if \texttt{Reason} is \texttt{false} the predicate \texttt{fail/0} is called as this will make \texttt{parallel/1} fail, and if \texttt{Reason} is of the form \texttt{exception(E)}, \texttt{E} is rethrown.

Note the use of an anonymous variable in the second argument of the patterns matching \texttt{false} and \texttt{exception(E)}. Using \texttt{Pid} here would work too, but might delay the failure of the call or the rethrowing of an error.

As soon as failure or error is detected, but before failing the entire call or rethrowing the error, the actor running \texttt{parallel/1} has a bit of tidying up to do, and this will be performed by a call to \texttt{tidy\_up\_all/1}. For this, it requires the complete list of worker actor pids, which has been passed along since its computation in the first call to \texttt{maplist/3} within the \texttt{parallel/1} definition.

The predicate \texttt{tidy\_up\_all/1} can be implemented as follows:  

\begin{verbatim}
tidy_up_all(Pids) :-
    maplist(tidy_up, Pids).

tidy_up(Pid) :-
    demonitor(Pid),
    exit(Pid, kill),
    drain_mailbox(Pid).

drain_mailbox(Pid) :-
    receive({
        Pid-_ ->
            drain_mailbox(Pid) ;
        down(_, Pid, _) ->
            drain_mailbox(Pid)
    },[
        timeout(0)
    ]).
\end{verbatim}

\noindent Here, one actor process at a time is killed by a call to \texttt{exit/2} (and recall that even if it is dead already, no error is raised). To avoid that the death of the process generates an additional \texttt{down} message, the monitoring of it must first be turned off. Finally, messages from the actor \texttt{Pid} that may have reached the mailbox during the execution of \texttt{parallel/1} are removed by a call to \texttt{drain\_mailbox/1}.

It is time to test our program and make sure that we get a speedup and that our four requirements are satisfied. In the following call to \texttt{parallel/1}, we simulate goals with long running times by combining simple unifications with calls to \texttt{sleep/1}. 

\begin{verbatim}
?- _Goals = [(X=a,sleep(1)),(Y=b,sleep(3)),(Z=c,sleep(2))],
   time(parallel(_Goals)).
 % 189 inferences, 0.000 CPU in 3.006 seconds
X = a,
Y = b,
Z = c.
?-
\end{verbatim}

\noindent Here the call was done in just 3 seconds rather than the 6 seconds it would take if running them sequentially. 

If one of the goals fail the entire call fails immediately, just as required by 3):

\begin{verbatim}    
?- _Goals = [(X=a,sleep(1)),(Y=b,fail),(Z=c,sleep(2))],
   time(parallel(_Goals)).
 % 105 inferences, 0.000 CPU in 0.001 seconds  
false.
?-
\end{verbatim}

\noindent And finally, as required by 4), if one of the goals is bad an error is thrown that can be caught:

\begin{verbatim}   
?- _Goals = [(X=a,sleep(1)),(Y=b,sleep(a)),(Z=c,sleep(2))],
   time(catch(parallel(_Goals),E,true)).
 % 126 inferences, 0.000 CPU in 0.001 seconds 
E = error(type_error(float, a), context(system:sleep/1, _)).
?-
\end{verbatim}

\noindent One thing to keep in mind is that if the goals in the list execute very quickly, the overhead of running \texttt{parallel/1} is likely to reduce any gains that might be had by running them in parallel. To make any noticeable difference, the predicate must be fed goals that run for a considerable time.

A predicate such as \texttt{parallel/1} is a valuable addition to the set of built-in predicates supported by Web Prolog, and it is encouraging to observe that it can be implemented with such succinctness. Naturally, there are multiple ways to implement this predicate. For example, consider Jan Wielemaker's implementation of \texttt{concurrent/3} (in his \texttt{library(thread)}\footnote{\url{https://www.swi-prolog.org/pldoc/doc/\_SWI\_/library/thread.pl?format\_comments=true\&show=src\#concurrent/3}}) which exhibits semantics sufficiently close to \texttt{parallel/1}. We are particularly encouraged by the fact that our implementation appears to be both simpler and more comprehensible than Wielemaker's. Code simplicity and clarity are important considerations, and code size matters.

\section{The first\_solution/2 meta-predicate}

Another example of parallel processing supported by Wielemaker, also available in \texttt{library(thread)}, is \texttt{first\_solution/2-3}. It tries alternative solvers in parallel, and as soon as one of the alternatives is successful, all the others are terminated.

Borrowing an example from Wielemaker, if it is unclear whether it is better to search a particular graph breadth-first or depth-first, we can use:

\begin{verbatim}
search_graph(Graph, Path) :-
    first_solution(Path, [ breadth_first(Graph, Path),
                           depth_first(Graph, Path)
                         ],
                   []).
\end{verbatim}

\noindent Applying roughly the same pattern as in the implementation of \texttt{parallel/1}, we implement \texttt{first\_solution/2} in Web Prolog below (leaving out the options that Wielemaker's version support):

\begin{verbatim}
first_solution(Solution, Goals) :-
    maplist(first_solve(Solution), Goals, Pids),
    wait_first(Pids, Solution).

first_solve(Solution, Goal, Pid) :-
    self(Self),
    spawn((call(Goal), Self ! Pid-Solution), Pid, [
        monitor(true)
    ]).

wait_first([], _) :- !, fail.
wait_first(Pids, Solution) :-
    receive({
        _ - Solution ->
            tidy_up_all(Pids) ;
        down(_, Pid, false) ->
            select(Pid, Pids, Rest),
            wait_first(Rest, Solution) ;
        down(_, _, exception(Error)) ->
            tidy_up_all(Pids),
            throw(Error)        
    }).
\end{verbatim}

\noindent Note the reuse of \texttt{tidy\_up\_all/1} from the implementation of \texttt{parallel/1}.

Using a simple simulation of two long-running goals, here is a trial run that shows that our predicate works as intended:

\begin{verbatim}
?- time(first_solution(X, [(sleep(2),X=a),(sleep(1),X=b)])).
 % 164 inferences, 0.000 CPU in 1.006 seconds
X = b.
?-
\end{verbatim} 


\section{Summary}

Taken together, the examples in this chapter illustrate how Web Prolog can
factor recurring concurrency and coordination patterns into reusable
behaviours. Erlang-style behaviours such as the server, supervisor, and
toplevel wrap actors in explicit protocols and runtime scaffolding for state
management, request--response handling, and fault tolerance, without binding
the caller's variables. Prolog-style generics such as \texttt{parallel/1}
and \texttt{first\_solution/2} show how the same actor substrate can be
used to implement higher-level meta-predicates that look and feel like
ordinary Prolog, yet exploit parallelism and robust error propagation
internally.

A single observation ties these examples together. In every case, the
application-specific logic --  the callback predicate that defines a
server's domain semantics, the goals passed to \texttt{parallel/1}, the
alternative solvers handed to \texttt{first\_solution/2} --  is written in
ordinary Prolog: pure predicates, relational specifications, familiar
recursive definitions. The concurrency, the fault handling, the protocol
machinery, and the cleanup discipline are supplied by the generic framework.
The developer writes the part that requires domain knowledge; the behaviour
provides the rest. This separation is not merely a convenience. It means
that the large existing body of Prolog programs and programming techniques
does not need to be abandoned or rewritten in order to participate in a
concurrent and distributed setting. It can be reused directly, wrapped in
the appropriate behaviour, and deployed as a networked service.

But the chapter also reveals a less obvious point. The Erlang-style
behaviours --  server, supervisor, toplevel --  have direct counterparts in
Erlang's OTP, and Web Prolog's versions are broadly comparable in structure
and length. The Prolog-style generics tell a different story. The
implementation of \texttt{parallel/1} requires roughly thirty lines of
code; \texttt{first\_solution/2} requires fewer than twenty. Both are
shorter and, we believe, more legible than their Erlang equivalents. The
reason is not a single feature but a combination: \texttt{maplist} handles
the orchestration loop, unification-based selective \texttt{receive}
correlates replies without manual bookkeeping, goals-as-terms pass work
units to spawned actors without serialisation or callback registration, and
\texttt{fail} triggers backtracking cleanup naturally. These are ordinary
Prolog mechanisms, none of them designed for concurrency, yet they compose
fluently with the actor primitives to yield concurrency abstractions that
would require considerably more scaffolding in a language without them. If
the Erlang-style behaviours show that Web Prolog can match OTP's
repertoire, the Prolog-style generics show where it can surpass it --  and
that difference is the clearest evidence this chapter offers for why the
combination of logic programming and actor programming is more than the sum
of its parts.